\documentclass[11pt]{article}
\usepackage[margin=1in]{geometry}
\usepackage{amsmath,amssymb,amsthm}
\usepackage[hidelinks]{hyperref}
\newtheorem{theorem}{Theorem}
\newtheorem{lemma}[theorem]{Lemma}
\newcommand{\mr}{\operatorname{mr}}
\newcommand{\F}{\mathbb F}
\title{Sharp cover certificates beyond finite Coxeter quotients}
\author{Internal research working proof}
\date{17 November 2026}
\begin{document}
\maketitle

The forty-eight ternary affine certificates prove meridional rank
four where no finite Coxeter reflection quotient gives that bound.
Here the negative assertion has a complete small replay, independent
of historical generating-set lists and workbook non-hits.
The separate universal reflection obstruction now also excludes
sharp infinite Coxeter reflection quotients on all sixteen alternating
cases. For the original eight, maximal Coxeter reflection rank is
exactly three by
explicit triangle quotients in
\path{REFLECTIVE_AUGMENTATION_RANK_GAP.tex}.
For the other thirty-two cases, infinite targets remain a separate
question. The computations and proofs below are internal,
unformalized and externally unreviewed; global priority remains open.

\section{The finite targets}

Let $\mathcal S$ consist of $13a_j$ for
\[
 j=373,375,391,392,3264,3266,3275,3276
\]
and $13n_j$ for
\[
\begin{split}
 j={}&1005,1006,1007,1011,1012,1013,1043,1044,\\
     &1045,1046,1047,1048,3842,3844,3855,3856.
\end{split}
\]
Add the twenty-four knots in $\mathcal K$ from
\path{A6_AFFINE_SHARP_FOUR.tex} to this set. Thus $|\mathcal S|=48$.
Each literal diagram has a retained four-seed propagation certificate.
The direct finite affine theorems prove $\mr(K)=b(K)=4$ throughout
this exact set; the original twenty-four also have seven-sheet proofs.

\begin{theorem}\label{thm:finite}
For every $K\in\mathcal S$ there is no epimorphism from its knot group
onto a finite Coxeter group of Coxeter rank four sending meridians to
Coxeter reflections. Consequently every finite Coxeter reflection
lower bound for these knots is at most three, whereas the retained
branched-cover bound is four.
\end{theorem}

The word \emph{reflection} is part of the statement. Arbitrary involutions
in these finite groups are not covered by the exclusion.

\begin{lemma}\label{lem:targets}
The only rank-four finite Coxeter systems which can receive such a
knot-group epimorphism are $A_4$, $D_4$ and $H_4$.
Moreover, a $D_4$ reflection epimorphism would give an $S_3$
transposition epimorphism.
\end{lemma}
\begin{proof}
In a Coxeter presentation, abelianization identifies generators
precisely along odd-labelled edges. Its abelianization is therefore
one copy of $C_2$ per component of that odd-edge graph. A knot group
has cyclic abelianization, so the odd-edge graph of a nontrivial
Coxeter quotient must be connected. This excludes reducible systems.
The finite irreducible rank-four classification leaves
$A_4,B_4,D_4,F_4,H_4$; see the classification discussion in
Blair--Kjuchukova--Morrison \cite[Section 2.1 and Section 3.1.1]{bkm}.
For $B_4$ and $F_4$ the even edge disconnects the odd-edge graph.
Their abelianization is $C_2^2$, so neither can occur.

In the signed-permutation model of $D_4$, reflections are signed
transpositions. The quotient $D_4\to S_4$ sends them to transpositions.
The further quotient $S_4\to S_4/V_4=S_3$ again sends every
transposition to a transposition. Surjectivity persists under
composition. The $A_4$ case is $S_5$ with its transpositions.
\end{proof}

\section{Complete transposition searches}

A fixed four-seed propagation plan determines every arc label from
its four seed labels. For an $S_r$ transposition coloring, conjugating
the target permits the first seed to be $(01)$. We exhaust all
\[
 \binom r2^3=27,\ 216,\ 1000
 \qquad (r=3,4,5)
\]
ordered tails, including repeated labels. Every crossing is checked
after propagation. On every one of the twenty-four literal PDs, in
each of these three degrees, the sole accepted normalized coloring
is the constant one. In particular there is no surjective coloring
onto $S_3$, $S_4$ or $S_5$. This excludes $A_4$ and $D_4$ by
Lemma~\ref{lem:targets}.

The propagation plan is only an enumeration reduction. Knot identity,
planarity, the one-component condition and every crossing relation
are separately checked. A different meridional generating family
cannot evade the enumeration: any knot-group homomorphism still
assigns labels to these diagram meridians, and the four seeds
generate them all.

\section{An exact sixty-reflection model for H4}

Put $\phi=(1+\sqrt5)/2$. Use twice unit normals $\mathbf U\in
\mathbb Z[\phi]^4$. The signed root set consists of
\[
\begin{array}{ll}
(\pm2,0,0,0)&\text{and all coordinate positions},\\
(\pm1,\pm1,\pm1,\pm1),&\\
(0,\pm1,\pm\phi,\pm(\phi-1))&\text{and all even coordinate permutations}.
\end{array}
\]
There are $8+16+96=120$ vectors, all with squared norm four.
Opposite vectors define the same reflection, leaving sixty labels.
Conjugation of reflections acts on normals by
\[
 \mathbf V\longmapsto
 \mathbf V-\tfrac12(\mathbf U\cdot\mathbf V)\mathbf U.
\]
The coefficient lies in $\mathbb Z[\phi]$. The checker verifies this
integrality, closure on all roots and involutivity exactly, using
$\phi^2=\phi+1$.

For completeness the model supplies a representation of the
abstract Coxeter presentation $H_4$ directly. Take
\[
\begin{split}
 \mathbf U_1&=(-2,0,0,0),\\
 \mathbf U_2&=(\phi,-1,\phi-1,0),\\
 \mathbf U_3&=(0,\phi,1-\phi,-1),\\
 \mathbf U_4&=(0,0,0,2).
\end{split}
\]
Their Gram matrix divided by two has diagonal two and adjacent
entries $-\phi,-1,-1$, with all other off-diagonal entries zero.
The induced exact permutations of the 120 signed roots have product
orders five, three, three on consecutive pairs, and two otherwise.
Thus the defining Coxeter relations hold. The normals are linearly
independent, so the four generating reflections have zero common
fixed space. Their conjugation action reaches all sixty root lines.

No faithfulness or computed group order is needed for the exclusion.
Composing a hypothetical epimorphism onto $H_4$ with this explicit
representation must generate those four reflections. All reflection
images are among the sixty root lines, and they are conjugate there.
The first seed can consequently be fixed to a single root label.

The checker exhausts all $60^3=216000$ ordered tails on each knot.
Among accepted colorings, the normalized counts are 121 on twelve
knots, 25 on $13n_{1013},13n_{1045}$, and one on the remaining ten.
Every accepted seed quartet has determinant zero in $\mathbb Z[\phi]$.
The reflected arc normals remain in its span. Hence its generated
reflection group fixes a nonzero vector in $\mathbb R^4$, and cannot
be the full explicit image, whose common fixed space is zero.
This excludes the final target and proves Theorem~\ref{thm:finite}.

\section{Verification and limits}

Run \texttt{verify\_finite\_coxeter\_exclusion.py} with standard Python.
It constructs the exact root model, verifies its Coxeter relations
and full-rank positive control, and enumerates every seed tail.
The positive simple-root control also realizes the familiar pretzel
cycle quotient with weights $(5,3,3,2)$, as in the pretzel construction
of Baader--Blair--Kjuchukova \cite[Section 2]{bbk}.
The checker rejects an invalid seed plan and dependent roots in its
full-rank control.

A second implementation, \texttt{verify\_finite\_coxeter\_modular.py},
rebuilds arc components and propagation in a different order and
uses arithmetic at the split primes eleven and thirty-one.
A scaled-root quartet determinant and its Galois conjugate each
have absolute value at most sixteen by Hadamard's inequality.
Their algebraic norm is an integer of absolute value at most 256.
Vanishing at a prime over eleven and a prime over thirty-one forces
that integer to be divisible by $341>256$, hence zero. This gives
an independent exact check of the determinant conclusion.
The second replay also checks root descriptors and the positive
control; its receipt states the precise modular-to-integral bridge.

These are complete finite searches on the retained inputs.
They do not exclude rank-four infinite Coxeter quotients,
rational-tangle extensions of known constructions, or every
previously resolved knot family. No global novelty inference follows
merely from finite Coxeter nonexistence.

\begin{thebibliography}{9}
\bibitem{bkm} R. Blair, A. Kjuchukova and N. Morrison,
Coxeter quotients of knot groups through 16 crossings,
arXiv:2208.09032v1, Sections 2.1 and 3.1.1.
\url{https://arxiv.org/abs/2208.09032v1}.
\bibitem{bbk} S. Baader, R. Blair and A. Kjuchukova,
Coxeter groups and meridional rank of links,
Math. Ann. 379 (2021), 1533--1551, Section 2.
\url{https://doi.org/10.1007/s00208-020-02124-z}.
\end{thebibliography}
\end{document}
